\section{Anexo: entregables del análisis sobre la inclusión del componente virtual}
\label{sec:anexo-virtualidad}

\begin{center}
{\large\bfseries\color{Azul-UCR} Propuesta de actualización del plan de estudios}\\[4pt]
{\bfseries para la inclusión del componente virtual}\\[10pt]
\begin{tabular}{@{}ll@{}}
\textbf{Nombre de la carrera:} & Bachillerato y Licenciatura en Matemática \\
\textbf{Unidad académica:} & Escuela de Matemática, Facultad de Ciencias \\
\textbf{Año:} & 2026 \\
\end{tabular}
\end{center}

\vspace{8pt}

Este documento sistematiza el proceso de análisis y reflexión realizado por la carrera para enriquecer el modelo de formación presencial mediante la inclusión del componente virtual en el plan de estudios, según las entregas estipuladas en la normativa propuesta para ello por la Vicerrectoría de Docencia. El análisis se articula con la reflexión pedagógica recogida en la sección \emph{Sobre la virtualidad} del capítulo de elementos pedagógicos y programas de los cursos.

\subsection{Entregable 1. Concepto de educación virtual de la carrera}

La educación virtual, para la carrera de Matemática, es un proceso formativo mediado por tecnologías digitales y por el entorno institucional de Mediación Virtual, que enriquece la interacción, el diálogo y la mediación pedagógica \emph{sin sustituir}, en la generalidad de los cursos, la presencialidad que exige la disciplina.

En Matemática, el aprendizaje se construye de manera central mediante la demostración, la escritura simbólica en vivo, la discusión de argumentos y la retroalimentación inmediata de la persona docente. Por ello, la virtualidad se concibe de dos modos, claramente diferenciados:

\begin{itemize}
    \item \textbf{Como apoyo a la presencialidad.} Uso de Mediación Virtual para organizar el curso, publicar materiales, recibir entregas, mantener comunicación asincrónica y disponer de recursos didácticos (notas, grabaciones complementarias, visualizaciones, repositorios de código). Esta mediación no reemplaza las horas de aula.
    \item \textbf{Como modalidad de las actividades lectivas}, parcial o total, \textbf{únicamente} en cursos prácticos con estilo de laboratorio computacional, en los que el trabajo en vivo con el computador es el medio natural de aprendizaje y puede ofrecer ventajas (implementación de algoritmos, experimentación numérica, visualización, control de versiones, elaboración de reportes reproducibles).
\end{itemize}

Esta concepción se sostiene en la evidencia y en la experiencia docente de la Escuela. El rendimiento académico \emph{inmediato} no siempre difiere de forma significativa entre la modalidad virtual y la presencial \cite{hachey2022post}; no obstante, la experiencia formativa sí se empobrece cuando se pierde la interacción directa. El estudiantado tiende a percibir la tecnología como un apoyo que no reemplaza el contacto con la persona docente \cite{Alshehri31122023}. En ciencias y matemática se ha documentado una preferencia por la presencialidad, con dificultades específicas de comprensión, de interacción, técnicas y de representación de conceptos \parencite{YorkovskyLevenberg2022}.

En consecuencia, la carrera no concibe un plan de estudios virtual ni bimodal de manera general. Concibe un plan \textbf{mayoritariamente presencial}, con mediación tecnológica de apoyo y con virtualidad lectiva restringida a laboratorios computacionales.

\subsection{Entregable 2. Rol de la educación virtual y perfiles en ambientes virtuales}

\subsubsection*{Rol de la educación virtual en la carrera}

La educación virtual cumple un rol \textbf{complementario y selectivo}:

\begin{enumerate}
    \item Complementa la clase presencial (organización, materiales, comunicación, entregas y, cuando proceda, grabaciones de apoyo).
    \item Permite, en cursos de laboratorio computacional, trasladar a entornos virtuales sincrónicos o asincrónicos aquellas actividades cuyo objeto es el uso de software científico, de modo que cada estudiante trabaje en su propio entorno de programación con acompañamiento docente.
    \item No se utiliza como sustituto de la construcción en vivo de demostraciones ni de la discusión continua de argumentos, propias de los cursos teóricos, en particular del énfasis en matemática pura y de los primeros años de la carrera.
\end{enumerate}

\subsubsection*{Perfil de la persona docente en ambientes virtuales}

La persona docente mantiene el rol de guía académica descrito en el marco pedagógico de la carrera (conocimiento disciplinar, gestión de la clase, clima matemático y retroalimentación). En los ambientes virtuales se espera, además, que:

\begin{itemize}
    \item Facilite, oriente y acompañe el trabajo en Mediación Virtual, con consignas explícitas y retroalimentación oportuna.
    \item Distinga con claridad qué actividades exigen presencialidad (demostración, discusión de pruebas, evaluación oral de argumentos) y cuáles pueden mediarse (laboratorios de código, visualización, entrega de reportes).
    \item Promueva la construcción colaborativa, la autogestión y la reflexión crítica, sin diluir el rigor ni la exigencia de escritura matemática precisa.
    \item En actividades asincrónicas, asuma un rol de tutoría: presencia visible en el entorno, plazos realistas y canales oficiales de consulta.
    \item Cumpla la formación en competencias digitales para la docencia dispuesta por la Vicerrectoría de Docencia (circular VD-15-2023) cuando el curso se catalogue con componente virtual según METICS.
\end{itemize}

\subsubsection*{Perfil de la persona estudiante en ambientes virtuales}

Se espera una persona estudiante activa, responsable y dispuesta al trabajo colaborativo, coherente con la conceptualización de la población estudiantil de la carrera. En ambientes virtuales, además:

\begin{itemize}
    \item Gestiona su tiempo, consulta los materiales publicados y entrega las actividades en los plazos establecidos.
    \item Participa en las sesiones sincrónicas cuando las haya (cámara y medios de escritura cuando la actividad lo requiera) y utiliza los espacios asincrónicos para preguntas y trabajo entre pares.
    \item Relaciona los nuevos aprendizajes con la experiencia de los cursos presenciales, de modo que la mediación tecnológica refuerce, y no fragmente, la formación matemática.
    \item Hace un uso ético del software, de los datos y de las herramientas de inteligencia artificial, cuando estas se autoricen, y documenta su propio trabajo (código, razonamiento, fuentes).
\end{itemize}

\subsection{Entregable 3. Mapeo de saberes}

\subsubsection*{Criterio general}

A partir de la experiencia de la pandemia de COVID-19 y de la reflexión pedagógica posterior, la carrera proyecta que los saberes asociados a la demostración, al formalismo, al pensamiento abstracto y a la discusión de argumentos se estimulan de manera presencial. La naturaleza de la disciplina dificulta trasladar a entornos virtuales la escritura en vivo, la intervención sobre un argumento y la retroalimentación inmediata que caracterizan una clase de matemática.

El mapeo no replica, por tanto, una virtualización del plan completo. Distingue:

\begin{enumerate}
    \item \textbf{Cursos teóricos (la mayoría del plan).} Permanecen presenciales. El uso de Mediación Virtual como apoyo a la docencia se cataloga como \textbf{bajo virtual}, según las disposiciones de METICS.
    \item \textbf{Cursos prácticos de laboratorio computacional.} Pueden ofrecerse en modalidad \textbf{parcialmente virtual} o, cuando el trabajo con el computador en vivo sea predominante, \textbf{virtual}. La modalidad concreta de cada oferta se declara en el programa y en la guía de horarios.
\end{enumerate}

\subsubsection*{Clasificación de cursos según la modalidad}

La Tabla~\ref{tab:virtualidad-cursos} resume la clasificación. Las siglas P, PV y V corresponden, respectivamente, a presencial, parcialmente virtual y virtual, de acuerdo con el uso ya adoptado en los programas de curso.

\begin{table}[H]
\centering
\small
\setlength{\tabcolsep}{4pt}
\renewcommand{\arraystretch}{1.15}
\begin{tabular}{|p{3.1cm}|p{5.4cm}|p{5.5cm}|}
\hline
\rowcolor[HTML]{00C0F3}
\textbf{Modalidad} & \textbf{Criterio pedagógico} & \textbf{Cursos (indicativos)} \\ \hline
Presencial, con apoyo en Mediación Virtual (bajo virtual METICS)
& Demostración en vivo, discusión de argumentos, escritura simbólica y evaluación que requiere retroalimentación inmediata.
& Bloque común teórico (MA-0151, MA-0152, MA-0251, MA-0252, MA-0351, MA-0361, MA-0451, MA-0461, entre otros); cursos del énfasis en matemática pura; seminarios y prácticas profesionales; optativas teóricas. \\ \hline
Parcialmente virtual (PV)
& Combinación de fundamentos teóricos con laboratorio de software, experimentación numérica o manejo de datos.
& MA-0341 Matemática Computacional; MA-0501 Análisis Numérico I; MA-0647 Modelación Matemática; MA-0920 EDP numéricas; CA-0204 Herramientas de Ciencia de Datos I; CA-0303 Estadística I; optativas de cómputo científico, análisis de datos y aprendizaje automático (MA-0795, MA-0798, MA-0836), según la oferta. \\ \hline
Virtual (V)
& El objeto de la actividad lectiva es el trabajo en vivo con el computador (implementación, análisis de datos, reportes reproducibles), de modo que la mediación virtual puede ser ventajosa.
& CA-0411 Análisis de Datos I (PV o V, según oferta). Otras ofertas de laboratorio computacional podrán declararse V cuando la persona docente lo justifique y se cumpla la normativa institucional. \\ \hline
\end{tabular}
\caption{Clasificación de cursos según el componente virtual. La lista del grupo PV/V es indicativa: la modalidad de cada oferta se fija en el programa del curso.}
\label{tab:virtualidad-cursos}
\end{table}

Los programas de los cursos de Ciencias Actuariales que toma el énfasis en matemática aplicada ya registran combinaciones P/PV o PV/V (CA-0204, CA-0303, CA-0411). En los cursos propios de Matemática con estilo de laboratorio, la virtualidad lectiva se admite en los términos de esta tabla, aun cuando la oferta ordinaria se mantenga presencial.

\subsubsection*{Mapeo de saberes del perfil y mediación virtual}

La Tabla~\ref{tab:mapeo-saberes-virtual} relaciona saberes del perfil de egreso con la posibilidad de desarrollarlos en entornos virtuales. Solo se detallan los saberes en los que la mediación tecnológica o el laboratorio computacional aportan de manera específica; los restantes saberes del perfil (demostrar, formalizar, comunicar oralmente un argumento, etc.) se consideran de desarrollo presencial.

\begin{landscape}
\begin{table}[H]
\centering
\small
\setlength{\tabcolsep}{4pt}
\renewcommand{\arraystretch}{1.15}
\begin{tabular}{|p{2.4cm}|p{5.2cm}|p{3.2cm}|p{4.4cm}|p{5.4cm}|}
\hline
\rowcolor[HTML]{00C0F3}
\textbf{Saber} & \textbf{Descripción breve} & \textbf{¿Virtual?} & \textbf{Cursos de anclaje} & \textbf{Estrategias de mediación} \\ \hline
SH03 & Experimenta con nuevas tecnologías para validar y conjeturar resultados.
& Parcial / sí en laboratorio
& MA-0341, MA-0501, MA-0795
& Guías de experimentación en software; bitácoras de conjeturas; foros de resultados en Mediación Virtual. \\ \hline
SH10 & Utiliza software científico para elaborar documentos, crear algoritmos, obtener soluciones y divulgar resultados.
& Sí (laboratorio)
& MA-0341, MA-0501, CA-0204, CA-0411, MA-0920
& Laboratorios sincrónicos o asincrónicos de código; repositorios; reportes en \LaTeX{} o Markdown. \\ \hline
SH13 & Usa la tecnología para explorar científicamente en áreas del conocimiento matemático y afines.
& Parcial / sí en laboratorio
& MA-0341, MA-0647, MA-0798, MA-0836
& Exploración dirigida con software; proyectos cortos; visualizaciones. \\ \hline
SH09 & Plantea problemas reales en términos matemáticos, mediante modelación, aproximación, simulación y pronóstico.
& Parcial
& MA-0647, MA-0501, CA-0411
& Ciclo de modelación con simulación numérica; la formulación y la interpretación se discuten de manera presencial. \\ \hline
SC13 & Conocimiento básico de la estructura de los sistemas numéricos (análisis numérico) y sus aplicaciones.
& Parcial
& MA-0501, MA-0920
& Implementación y prueba de algoritmos en laboratorio; la teoría de convergencia y estabilidad se desarrolla de manera presencial. \\ \hline
SH05, SH08 y SH12 & Leer, escribir y pensar de manera abstracta y formal en lenguaje matemático.
& No como actividad lectiva principal
& Cursos teóricos del plan
& Mediación Virtual solo como repositorio y canal de entrega; la mediación pedagógica es presencial. \\ \hline
\end{tabular}
\caption{Mapeo de saberes del perfil de egreso respecto de la mediación virtual. Los códigos corresponden al perfil de la carrera.}
\label{tab:mapeo-saberes-virtual}
\end{table}
\end{landscape}

\subsubsection*{Corroboración de coherencia curricular}

De manera análoga a lo solicitado en la plantilla institucional, se corroboró que los objetivos de los programas respondan al perfil de egreso, que estén enfocados en el aprendizaje y sean evaluables, y que la carga académica sea acorde con el creditaje. Ese trabajo se recoge en los programas del Anexo de cursos y en el capítulo de estructura curricular. No se virtualizan saberes cuya evaluación exige la producción y la defensa presencial de un argumento matemático.

\begin{landscape}
\subsection{Entregable 4. Sistematización de las necesidades para la inclusión del componente virtual}

\begin{table}[H]
\centering
\small
\setlength{\tabcolsep}{4pt}
\renewcommand{\arraystretch}{1.15}
\begin{tabular}{|p{5.3cm}|p{5.5cm}|p{5.8cm}|p{4.8cm}|}
\hline
\rowcolor[HTML]{00C0F3}
\textbf{Requerimientos tecnológicos del estudiantado} &
\textbf{Necesidades de capacitación para el personal docente y administrativo} &
\textbf{Consideraciones para la interacción y el aprendizaje colaborativo y autónomo} &
\textbf{Otras necesidades detectadas} \\ \hline
Computadora con capacidad para ejecutar software científico (Python, Julia, R, SageMath u otros indicados en el programa), tableta u otro dispositivo electrónico.
&
Cumplir las 30 horas de Competencias Digitales para la Docencia, según la circular VD-15-2023, en particular el personal que imparta cursos catalogados con componente virtual.
&
Tomar en cuenta las motivaciones del estudiantado y sus necesidades en el uso de herramientas tecnológicas, sobre todo en el tránsito de la secundaria a los primeros cursos de laboratorio.
&
Dificultad técnica de la escritura matemática en vivo en entornos virtuales (pizarra digital, tableta con lápiz); se evita trasladar a virtualidad las demostraciones.
\\ \hline
Conexión a internet estable para las actividades sincrónicas de los cursos PV o V.
&
Capacitación continua en programación científica, software matemático y estadístico, preparación de documentos con \LaTeX{} y uso de Mediación Virtual, según el plan de desarrollo docente (METICS y Escuela de Matemática).
&
Incluir actividades de creación, exploración, observación, discusión y pensamiento crítico; en laboratorio, bitácoras de experimentación y proyectos cortos.
&
Riesgo de conexión deficiente o de desconexión por factores no controlables; las actividades sincrónicas virtuales deben prever alternativas asincrónicas equivalentes.
\\ \hline
Acceso al laboratorio de cómputo de la Escuela de Matemática o a equipos de las bibliotecas, para quien no disponga de equipo adecuado en su lugar de residencia.
&
Inducción al plan 2028 y elaboración de guías metodológicas que expliciten qué partes de cada curso son presenciales y cuáles pueden mediarse.
&
La persona docente mantiene un rol de acompañamiento o tutoría en las actividades asincrónicas, con retroalimentación sobre código y sobre el razonamiento.
&
Licencias, repositorios y entornos institucionales para software científico y para control de versiones, coordinados con la Escuela y con METICS.
\\ \hline
Cuenta institucional y dominio básico de Mediación Virtual.
&
Actualización disciplinar en ciencia de datos, aprendizaje automático, modelación y computación científica, para sostener las optativas con componente virtual.
&
Mantener comunicación constante por los medios oficiales (Mediación Virtual y correo institucional), sin sustituir la consulta presencial en horas de atención.
&
Evaluación: en cursos teóricos, las pruebas que juzgan la producción de argumentos se conservan presenciales; en laboratorio, se pueden evaluar entregas de código y reportes en el entorno virtual, con criterios de autoría y de rigor.
\\ \hline
\end{tabular}
\caption{Necesidades identificadas para la inclusión del componente virtual en el plan de estudios de Matemática.}
\label{tab:necesidades-virtualidad}
\end{table}
\end{landscape}

En síntesis, la carrera incorpora el componente virtual de manera \textbf{restringida y justificada}: como apoyo sistemático a la presencialidad en todo el plan, y como modalidad lectiva solo en cursos prácticos de laboratorio computacional, donde el uso del computador en vivo puede favorecer el aprendizaje. Con ello se busca una formación más integral, sin sacrificar la profundidad que aporta la interacción directa en la demostración matemática.
